%%%%%%%%%%%%%%%%%%%%%%%%%%%%%%%%%
%
%       EXERCÍCIO
%
%%%%%%%%%%%%%%%%%%%%%%%%%%%%%%%%%

\ifdefstring{\atividade}{prova}{%
    \ifdefstring{\modo}{objetivo}{%
        \renewcommand{\valorquestao}{\ValorQObj\ ponto}
    }{%
        \renewcommand{\valorquestao}{\ValorQDisc\ pontos}
    }%
}%

\begin{exercicioBanco}[\valorquestao]
Uma empresa produz camisetas sob encomenda. O custo total de produção é composto por um custo fixo, mais um custo variável proporcional ao número de camisetas produzidas. O custo total cresce de forma linear, e a tabela apresenta esse custo para algumas quantidades produzidas. 
%
\begin{center}
    \begin{tabular}{|C{200pt}|C{50pt}|C{50pt}|}
        \hline
        Quantidade de camisetas & 20 & 40 \\
        \hline
        Custo total (R\$) & 700{,}00 & 1100{,}00 \\
        \hline \hline
    \end{tabular}
\end{center}
%
Determine o custo total, em real, para a produção de 60 camisetas.

% Define as alternativas
\newcommand{\alternativas}{%
    \begin{center}
        \begin{tabularx}{\textwidth}{XXXXX}
            (a) R\$ \(1500{,}00\). &
            (b) R\$ \(1400{,}00\). &
            (c) R\$ \(1600{,}00\). &
            (d) R\$ \(1700{,}00\). &
            (e) R\$ \(1800{,}00\).
        \end{tabularx}
    \end{center}
}

% Resposta correta
\newcommand{\resposta}{A}

% Lógica condicional para exibição
\ifdefstring{\atividade}{lista}{%
    \alternativas
    \vspace{0.5em}
    
    \noindent\textbf{Resposta:} letra \textbf{\resposta}.
}{%
    \ifdefstring{\modo}{objetivo}{%
        \alternativas
    }{%
        % modo = discursiva → não mostra alternativas
    }
}
\end{exercicioBanco}

